Soit une fonction $f$ int\'egrable sur $(a,b)$ ($-\infty \leq a < b \leq +\infty$).
On veut calculer ici \`a l'aide de m\'ethodes num\'eriques :

\[
I = \int_a^b f(x) dx.
\]

En particulier, on va chercher \`a \'evaluer les int\'egrales suivantes :

\begin{align*}
I_1 =& \int_0^{2\pi} |\cos(x)| dx,\\
I_2 =& \int_0^1 \frac{dx}{1 + 4x^2}.
\end{align*}

\begin{enumerate}

\item SAGE est normalement capable de calculer symboliquement ({\it i.e.} analytiquement) des int\'egrales.
Essayer de calculer $I_1$ et $I_2$ ainsi. Que se passe-t'il pour $I_1$ ?\\

{\it Indication: commande $\verb+integral()+$.}

\item On va \'evaluer $I_1$ et $I_2$ num\'eriquement \`a l'aide d'une formule de quadrature compos\'ee.
Il faut pour cela d\'efinir une subdivision de l'intervalle $(a,b)$. Soit donc $n \geq 1$ et 

$$a = \alpha_0 < \alpha_1 < \ldots < \alpha_n = b$$

une subdivision de $(a,b)$. Choisir une valeur de $n$ (pas trop \'elev\'ee au d\'ebut) et construire
sous SAGE la liste des $\alpha_i$ (on pourra choisir une subdivision r\'eguli\`ere).

\item Calculer num\'eriquement $I_1$ \`a l'aide d'une formule de quadrature compos�e sur la subdivision 
d\'efinie pr\'ec\'edemment. On utilisera pour les quadratures \'el\'ementaires :

\begin{enumerate}
 \item la m\'ethode des rectangles \`a gauche.
 \item la m\'ethode du point milieu.
 \item la m\'ethode des trap\`ezes.
 \item la m\'ethode de Simpson.
\end{enumerate}

Quelle est l'erreur d'int\'egration $\varepsilon(I_1)$ pour chaque m\'ethode ?

\item Faire varier $n$ et tracer le graphique de l'erreur d'int\'egration $\varepsilon(I_1)$
en fonction de $n$ {\it pour chaque m\'ethode}. Quelle m\'ethode vous semble la meilleure ?
Pourquoi ?

\item Faire de m\^eme avec $I_2$. Quelles sont vos conclusions ?

\end{enumerate}

